\documentclass[10pt,a4paper]{article}

% ----------------- Paquetes -------------------%
\usepackage[T1]{fontenc}
\usepackage[utf8]{inputenc}
\usepackage[a4paper,margin=2.5cm]{geometry}
\usepackage{mathptmx}             
\usepackage{changepage} 
\usepackage{setspace}
\usepackage[spanish]{babel}
\usepackage[labelfont=bf,labelsep=space]{caption}
\usepackage{amssymb}
\usepackage{amsmath}
\usepackage{ebgaramond}
\usepackage{placeins}
\usepackage{etoolbox}
\usepackage{setspace}
\usepackage{parskip} 
\usepackage{tcolorbox}
\usepackage{needspace}
\usepackage{xparse}
\usepackage{ocgx2}
\usepackage{hyperref}
\usepackage{hypcap}
\usepackage{soul}\sethlcolor{yellow}% resaltado amarillo: \hl{texto} (Panel TCEE, ⌃H)



%%%%%%%%%%%%%%%%%%%%%%%%%%%%%%%%%%%%%%%%%%%%%%%%%%%%%%%%%%%%%%%%
% --- Fija primero tus valores globales "buenos" ---
\setstretch{1.2}                 % Interlineado global
\setlength{\parskip}{0.7\baselineskip}
\setlength{\parindent}{0pt}
\newlength{\GlobalParskip}
\setlength{\GlobalParskip}{\parskip}

\newcounter{eqnum}[subsection]
\renewcommand{\theeqnum}{\arabic{eqnum}}
\newcounter{alphaitem}
\newcounter{numitem}

%% ------------------- Recuadros----------------------%%
\newcommand{\puntosclave}[1]{%
  \begin{tcolorbox}[
    colframe=black,
    title={Puntos clave},
    before upper={\setlength{\parindent}{0.5cm}\setlength{\parskip}{0.5\baselineskip}}
  ]
  #1
  \end{tcolorbox}
}

\newcommand{\modificaciones}[1]{
  \begin{tcolorbox}[
    colback=white,
    colframe=red,
    title={Modificaciones a realizar},
    before upper={\setlength{\parindent}{0.5cm}\setlength{\parskip}{0.5\baselineskip}}
  ]
  #1
  \end{tcolorbox}
}

%% ------------------- Preguntas TEST----------------------%%
% Opciones: radio (single-choice)
\NewDocumentCommand{\OptRadio}{m m m}{%
  \noindent\CheckBox[radio,name=#1,value=#2,width=1.2ex,height=1.2ex]{}%
  \;\textbf{#2.}~#3\par
}

% Opciones: checkbox (multi-choice)
\NewDocumentCommand{\OptCheck}{m m}{%
  \noindent\CheckBox[name=#1,width=1.2ex,height=1.2ex,bordercolor={0 0 0}]{}%
  \;#2\par
}

% Pregunta single-choice (radio)
% Args: 1=id, 2=enunciado, 3=opciones, 4=solución+justificación, 5=imagen (o {}), 6=origen
\NewDocumentCommand{\PreguntaTestSingle}{m m m m m m}{%
  \par\medskip
  \noindent\textbf{#1.}~#2\par\smallskip

    % Imagen (si no es {})
    \IfBlankTF{#5}{}{%
      \begin{center}
        \includegraphics[width=0.75\linewidth]{#5}
      \end{center}
    }

    \begin{Form}
      #3
    \end{Form}

    \par\smallskip
    \switchocg{sol-#1}{\fbox{Mostrar / ocultar solución}}%
    \hfill{\footnotesize #6}\par\smallskip

    \begin{ocg}{Solución}{sol-#1}{0}
      \noindent #4
    \end{ocg}
  \par\medskip
}

% Pregunta multi-choice (checkboxes)
\NewDocumentCommand{\PreguntaTestMulti}{m m m m m m}{%
  \par\medskip
  \noindent\textbf{#1.}~#2\par\smallskip

    % Imagen (si no es {})
    \IfBlankTF{#5}{}{%
      \begin{center}
        \includegraphics[width=0.75\linewidth]{#5}
      \end{center}
    }
    \begin{Form}
      #3
    \end{Form}

    \par\smallskip
    \switchocg{sol-#1}{\fbox{Mostrar / ocultar solución}}%
    \hfill{\footnotesize #6}\par\smallskip

    \begin{ocg}{Solución}{sol-#1}{0}
      \noindent #4
    \end{ocg}
  \par\medskip
}

%% ------------------- Listas----------------------%%
    % --- Lista alfabética ---
    \newenvironment{la}
      {\begin{list}{\alph{alphaitem})}{%
          \usecounter{alphaitem}%
          \setlength{\leftmargin}{1cm}%
          \setlength{\parskip}{3pt}% 
          \setlength{\itemsep}{0pt}% ← espacio entre ítems
          \setlength{\parsep}{0pt}%  ← párrafos dentro de un ítem
      }}
      {\end{list}}
      
    % --- Lista numérica ---
    \newenvironment{lnum}
      {\begin{list}{\arabic{numitem}.}{%
          \usecounter{numitem}%
          \setlength{\leftmargin}{1cm}%
          \setlength{\parskip}{3pt}% 
          \setlength{\itemsep}{0pt}% ← espacio entre ítems
          \setlength{\parsep}{0pt}%  ← párrafos dentro de un ítem
      }}
      {\end{list}}

%% ------------------- Preguntas Test ----------------------%%
      \usepackage{xparse} % si ya lo tienes, no lo repitas

    \NewDocumentEnvironment{test}{m m o}{%
      \begin{tcolorbox}[
        colback=white,
        colframe=black!15,
        boxrule=0.6pt,
        arc=2mm,
        left=2mm,right=2mm,top=2mm,bottom=2mm
      ]
      \centering
      \includegraphics[width=\linewidth]{#1}
      \end{tcolorbox}
    
      \vspace{2mm}
    
      % Caja SOLO para texto (SÍ partible y mantiene márgenes al partir)
      \begin{tcolorbox}[
        enhanced,
        breakable,
        colback=black!3,
        colframe=black!25,
        boxrule=0.6pt,
        arc=2mm,
        left=2mm,right=2mm,top=1.5mm,bottom=1.5mm
      ]
      \textbf{Respuesta:} \textbf{#2}%
      \IfValueT{#3}{\par\smallskip \textbf{Justificación:} #3}
      \end{tcolorbox}
    }{}
      
% ------------------- Imágenes----------------------%
    \graphicspath{{Img/}}% Rutas por defecto 
    % --- Imagen adaptativa ---
    % Registros de dimensión definidos una sola vez
    \newdimen\imgcap
    \newdimen\imgmin
    \newdimen\imgmax
    \newdimen\imgremain
    \newdimen\imgh
    \newdimen\imgneed
    
    \newcommand{\imagenfit}[3]{%
      \addvspace{10pt}% separación antes
      \par\begingroup
        % Parámetros de composición (asignaciones, no \newdimen)
        \imgcap=1\baselineskip            % alto estimado de título+fuente
        \imgmin=0.15\textheight
        \imgmax=0.15\textheight
        \imgremain=\pagegoal
        \ifdim\pagegoal=\maxdimen \imgremain=0pt\fi
        \advance\imgremain by -\pagetotal
        \advance\imgremain by -\imgcap
        % Decisión de altura destino
        \ifdim\imgremain<\imgmin
          \newpage
          \imgh=0.20\textheight
        \else
          \imgh=\imgremain
          \ifdim\imgh>\imgmax \imgh=\imgmax \fi
        \fi
        % Reservar espacio para evitar cortes del bloque completo
        \imgneed=\imgh \advance\imgneed by \imgcap
        \Needspace{\imgneed}
        % Cuerpo
        \noindent\begin{minipage}{\textwidth}
          \centering
          \captionof{figure}{\textemdash\ #2}\par\vspace{0.3em}% título arriba
          \includegraphics[width=\linewidth,height=\imgh,keepaspectratio]{\detokenize{#1}}\par
          \vspace{0.3em}\small\textit{Fuente: }#3% fuente abajo
        \end{minipage}%
      \endgroup
      \par\addvspace{10pt}%
    }

%------------ Bloque de RECUADRO ESPECIAL -------------%
    \newenvironment{specialblock}[1]{%
      \begin{quote} % crea sangría a izquierda y derecha
      \small % texto más pequeño
      \noindent\textbf{#1}\\[-0.3em] % título en negrita
      \rule{\linewidth}{0.4pt}\\[0.6em]}{
      \end{quote}}
      
%-------------------------- Portada -----------------------%
\newcommand{\oldtitle}[1]{\def\OldTitle{#1}}
\newcommand{\changerationale}[1]{\def\ChangeWhy{#1}}

\newcommand{\changebox}{%
\begin{tcolorbox}[title={Historial de cambios del tema}]
\textbf{Título anterior:} \OldTitle\par
\textbf{Motivo del cambio:} \ChangeWhy
\end{tcolorbox}
}

%-------------------------- Author Font -----------------------%
    \newcommand{\authorfont}[1]{{\fontfamily{EBGaramond-TLF}\selectfont #1}}

%-------------------------- Anexos -----------------------%
    \makeatletter
    \newcounter{anexo}
    \renewcommand{\theanexo}{\Roman{anexo}}
    \newcommand{\anexo}[1]{%
      \clearpage
      \phantomsection
      \refstepcounter{anexo}%
      \noindent\textbf{Anexo \theanexo.- }#1\par\medskip
      \addcontentsline{stc}{section}{Anexo \theanexo.- #1}%
    }
    \makeatother

    %-------------------------- Lista de figuras (personal) -----------------------%
\makeatletter
\newcommand{\MyListOfFigures}{%
  \clearpage
  \phantomsection
  {\centering\bfseries\Large Índice de figuras\par}%
  \medskip
  % Entrada en nuestro índice de contenido (stc)
  \addcontentsline{stc}{section}{Índice de figuras}%
  % Recuperación del listado (sin cabecera propia)
  \begingroup
    \parindent=0pt\parskip=2pt
    \@starttoc{lof}%
  \endgroup
}
\makeatother

%-------------------------- Bibliografía -----------------------%
\makeatletter
% Contador para las referencias
\newcounter{myref}
% Cabecera de la bibliografía, entra en el índice 'stc'
\newcommand{\MyBibliography}{%
  \clearpage
  \phantomsection
  {\centering\bfseries\Large Bibliografía y referencias\par}%
  \medskip
  \addcontentsline{stc}{section}{Bibliografía y referencias}%
  \setcounter{myref}{0}%
}
\newcommand{\MyBibItem}[4]{%
  \refstepcounter{myref}%
  \noindent[\themyref]\ #1.\ \emph{#2}. #3. #4\par
}
\makeatother

%--------- Establecer cuatro niveles de índice --------%
    \setcounter{tocdepth}{4}   
    \setcounter{secnumdepth}{4}% numera hasta \paragraph

%%%%%%%%%%%%%%%%%%%%%%%%%%%%%%%%

\makeatletter

    %--------- Interlineado Global --------%
    \edef\GlobalBaseLineStretch{\baselinestretch}
    
    %--------- Espaciado de seccinones --------%
    \renewcommand\section{\@startsection{section}
      {1}{\z@}
      {2.5ex \@plus 1ex \@minus .2ex} % espacio antes
      {1.5ex \@plus .2ex}             % espacio después
      {\normalfont\Large\bfseries}    % formato del título
    }
    \renewcommand\subsection{\@startsection{subsection}{2}{\z@}
      {3ex \@plus 0.8ex \@minus .2ex}
      {1.5ex \@plus .2ex}
      {\normalfont\large\bfseries}
    }
    \renewcommand\subsubsection{\@startsection{subsubsection}{3}{\z@}
      {3ex \@plus 0.5ex \@minus .2ex}
      {1ex \@plus .2ex}
      {\normalfont\normalsize\bfseries}
    }
    \renewcommand\paragraph{\@startsection{paragraph}{4}{\z@}%
    {-2ex \@plus -0.5ex \@minus -.2ex}% espacio antes
    {0.8ex \@plus .2ex}% espacio después
    {\normalfont\normalsize\itshape}}% ← cursiva en lugar de negrita

    %--------- Ecuaciones --------%   
    % Variables globales para almacenar títulos y notas
    \gdef\leq@current@section{}%
    \gdef\leq@current@subsection{}%
    \gdef\leq@last@printed@section{}%
    \gdef\leq@last@printed@subsection{}%
    
    % Comando para añadir nota (invisible en el cuerpo)
    \newcommand{\eqnote}[1]{%
      \addcontentsline{leq}{leqnote}{#1}%
    }
    
    % Comando para ecuaciones
    \newcommand{\eqblock}[2]{%
      % Verificar si necesitamos imprimir el título de sección
      \ifx\leq@current@section\leq@last@printed@section\else
        \addcontentsline{leq}{leqsection}{\leq@current@section}%
        \global\let\leq@last@printed@section\leq@current@section
        \gdef\leq@last@printed@subsection{}% Reset subsección al cambiar sección
      \fi
      % Verificar si necesitamos imprimir el título de subsección
      \ifx\leq@current@subsection\leq@last@printed@subsection\else
        \ifx\leq@current@subsection\@empty\else
          \addcontentsline{leq}{leqsubsection}{\leq@current@subsection}%
          \global\let\leq@last@printed@subsection\leq@current@subsection
        \fi
      \fi
      % Ecuación
      \refstepcounter{eqnum}%
      \begin{equation*}
        #1\tag{E.\theeqnum}
      \end{equation*}%
      \addcontentsline{leq}{leqt}{[E.\theeqnum] -- #2}%
      \addcontentsline{leq}{leqm}{#1}%
    }
    
    %%% ----- Índice de ecuaciones ----------%%%
    % Tamaño para []
    \newcommand{\leqIndexFont}{\normalsize}
    \newcommand{\leqTitleFont}{\tiny}
    \newcommand{\leqNoteFont}{\tiny}
    % Cabecera de sección 
    \newcommand*\l@leqsection[2]{%
      \addvspace{25pt}% Mayor separación antes de nueva sección
      \noindent\leqIndexFont
      \makebox[\linewidth]{\rule{.38\linewidth}{0.4pt}\ \ \textbf{  \MakeUppercase{#1}}\ \ \rule{.38\linewidth}{0.4pt}}%
      \par\vspace{-10pt}
    }
    % Cabecera de subsección 
    \newcommand*\l@leqsubsection[2]{%
      \par\addvspace{15pt}%
      \noindent\leqIndexFont #1
      \par\vspace{-7pt}
    }
    % Título de ecuación 
    \newcommand*\l@leqt[2]{%
    \par\addvspace{10pt}%
      \par\noindent\leqTitleFont #1\par
    }
    % Miniatura de ecuación
    \newcommand*\l@leqm[2]{%
      \vspace{0.3\baselineskip}%
      \par\scriptsize\noindent\hspace*{1.5em}\(#1\)\par%
    }
    % Nota de ecuación 
    \newcommand*\l@leqnote[2]{%
      \vspace{0.2\baselineskip}%
      \par\noindent\leqNoteFont\hspace*{1.5em}\textbf{Nota.--} #1\par\vspace{0.5\baselineskip}%
    }
    % Generación del índice
    \newcommand{\mylistofequations}{%
      \clearpage
      \phantomsection
      % Título visual de la lista (ajústalo a tu gusto):
      {\centering\bfseries\Large Índice de ecuaciones\par}
      % Entrada en nuestro índice 'stc':
      \addcontentsline{stc}{section}{Índice de ecuaciones}%
      % Llamada a tu lista real (usa el nombre exacto de tu macro):
      \@starttoc{leq}%
    }
    
    %--------- Captura de títulos de sección --------%
    \let\leq@original@section\section
    \let\leq@original@subsection\subsection
    % Flag para saber si estamos en una sección asterisco
    \newif\ifLeqStarSection
    \LeqStarSectionfalse
    
    % Redefinir \section CON CONTROL DE ASTERISCO
    \renewcommand{\section}{%
      \@ifstar{\leq@section@star}{\@dblarg\leq@section@nostar}%
    }
    \def\leq@section@star#1{%
      \global\LeqStarSectiontrue
      \gdef\leq@current@section{#1}%
      \gdef\leq@current@subsection{}%
      \leq@original@section*{#1}%
      \global\LeqStarSectionfalse
    }
    \def\leq@section@nostar[#1]#2{%
      \global\LeqStarSectionfalse
      \gdef\leq@current@section{#2}%
      \gdef\leq@current@subsection{}%
      \leq@original@section[#1]{#2}%
    }
    % Redefinir \subsection
    \renewcommand{\subsection}{%
      \@ifstar{\leq@subsection@star}{\@dblarg\leq@subsection@nostar}%
    }
    \def\leq@subsection@star#1{%
      \global\LeqStarSectiontrue
      \gdef\leq@current@subsection{#1}%
      \leq@original@subsection*{#1}%
      \global\LeqStarSectionfalse
    }
    \def\leq@subsection@nostar[#1]#2{%
      \global\LeqStarSectionfalse
      \gdef\leq@current@subsection{#2}%
      \leq@original@subsection[#1]{#2}%
    }

    % ----- Restauración de valores Globales de estilo ----------%
    % Flags
    \newif\ifInFootnote
    \InFootnotefalse
    \pretocmd{\@footnotetext}{\InFootnotetrue}{}{}
    \apptocmd{\@footnotetext}{\InFootnotefalse}{}{}
    % Profundidad de listas (soporta anidadas)
    \newcount\MyListDepth \MyListDepth=0
    \AtBeginEnvironment{la}{\advance\MyListDepth by 1}
    \AfterEndEnvironment{la}{\advance\MyListDepth by -1}
    \AtBeginEnvironment{lnum}{\advance\MyListDepth by 1}
    \AfterEndEnvironment{lnum}{\advance\MyListDepth by -1}
    \AtBeginEnvironment{itemize}{\advance\MyListDepth by 1}
    \AfterEndEnvironment{itemize}{\advance\MyListDepth by -1}
    \AtBeginEnvironment{enumerate}{\advance\MyListDepth by 1}
    \AfterEndEnvironment{enumerate}{\advance\MyListDepth by -1}
    % Restauración diferida: hazlo al INICIO del próximo párrafo real
    \newif\if@pendingRestore
    \@pendingRestorefalse
    \newcommand{\RequestRestore}{\global\@pendingRestoretrue}
    \everypar{%
      \if@pendingRestore
        \global\@pendingRestorefalse
        \ifInFootnote\else
          \ifnum\MyListDepth=0
            \setlength{\parskip}{\GlobalParskip}%
            \setstretch{\GlobalBaseLineStretch}%
          \fi
        \fi
      \fi
    }
    % Al final de bloques:
    \AfterEndEnvironment{equation}{\RequestRestore}
    \AfterEndEnvironment{equation*}{\RequestRestore}
    \AfterEndEnvironment{la}{\RequestRestore}
    \AfterEndEnvironment{lnum}{\RequestRestore}
    % Al final de macros:
    \apptocmd{\eqblock}{\RequestRestore}{}{}
    \apptocmd{\imagenfit}{\RequestRestore}{}{}
    
\makeatother

% ===================== ToC propio con 4 niveles (único bloque) =====================
\makeatletter

% --- Escritores: añadimos una línea al .stc en cada cabecera numerada ---
% Guardamos las versiones actuales para llamar al comportamiento normal
\let\MyTOC@orig@section\section
\let\MyTOC@orig@subsection\subsection
\let\MyTOC@orig@subsubsection\subsubsection
\let\MyTOC@orig@paragraph\paragraph

% SECTION
\renewcommand{\section}{%
  \@ifstar{\MyTOC@section@star}{\@dblarg\MyTOC@section@nostar}%
}
\def\MyTOC@section@star#1{%
  \MyTOC@orig@section*{#1}% sin entrada al .stc
}
\def\MyTOC@section@nostar[#1]#2{%
  \MyTOC@orig@section[{#1}]{#2}% comportamiento normal (escribe al .toc)
  % Escribimos también nuestra línea al .stc (texto con \numberline)
  \addcontentsline{stc}{section}{\protect\numberline{\thesection}#2}%
}

% SUBSECTION
\renewcommand{\subsection}{%
  \@ifstar{\MyTOC@subsection@star}{\@dblarg\MyTOC@subsection@nostar}%
}
\def\MyTOC@subsection@star#1{%
  \MyTOC@orig@subsection*{#1}%
}
\def\MyTOC@subsection@nostar[#1]#2{%
  \MyTOC@orig@subsection[{#1}]{#2}%
  \addcontentsline{stc}{subsection}{\protect\numberline{\thesubsection}#2}%
}

% SUBSUBSECTION
\renewcommand{\subsubsection}{%
  \@ifstar{\MyTOC@subsubsection@star}{\@dblarg\MyTOC@subsubsection@nostar}%
}
\def\MyTOC@subsubsection@star#1{%
  \MyTOC@orig@subsubsection*{#1}%
}
\def\MyTOC@subsubsection@nostar[#1]#2{%
  \MyTOC@orig@subsubsection[{#1}]{#2}%
  \addcontentsline{stc}{subsubsection}{\protect\numberline{\thesubsubsection}#2}%
}

% PARAGRAPH
\renewcommand{\paragraph}{%
  \@ifstar{\MyTOC@paragraph@star}{\@dblarg\MyTOC@paragraph@nostar}%
}
\def\MyTOC@paragraph@star#1{%
  \MyTOC@orig@paragraph*{#1}%
}
\def\MyTOC@paragraph@nostar[#1]#2{%
  \MyTOC@orig@paragraph[{#1}]{#2}%
  % Si no numeras los \paragraph, cambia \theparagraph por texto plano sin \numberline
  \addcontentsline{stc}{paragraph}{\protect\numberline{\theparagraph}#2}%
}

% --- Impresor del índice propio (lee el .stc) ---
\newcommand{\MyTOC}{%
  \par\bigskip
  {\centering\bfseries\Large Índice de contenido\par}%
  \medskip
  \begingroup
    \parindent=0pt\parskip=2pt
    \@starttoc{stc}%
  \endgroup
}

\makeatother
% ===================== fin ToC propio =====================


%%%%%%%%%%%%%%


% --- Datos del tema ---
\title{Tema 3.A.9 -- Teoría de la demanda del consumidor (II)\\
\large Dualidad e integrabilidad de las preferencias. Sistemas de demanda utilizados en estudios empíricos. Medidas de cambio en el binestar}
\author{Marcelo Ortega}
\date{Febrero 2026}
\newcommand{\TemaCode}{3A09}

\oldtitle{Tema 3.A.8. La dualidad en la teoría de la demanda del consumidor y sus aplicaciones}
\changerationale{
    Se explicitan los contenidos necesarios del tema para guiar al opositor y garantizar un mínimo común. Se añade una referencia a aplicaciones empíricas de sistemas de demanda, siguiendo la línea iniciada por Angus Deaton.
}

\begin{document}


\maketitle
\changebox

\newpage

% --- Página de resumen y modificaciones ---
\puntosclave{ %% Escribir puntos clave Apuntes CECO
    }
\vspace{1cm}

\modificaciones{
    
    }

\newpage
\MyTOC
\newpage

\mylistofequations
\clearpage

\MyListOfFigures
\clearpage


\pagenumbering{arabic}
\setcounter{page}{1}


\section{Introducción}\label{intro}
ALFRED MARSHALL, en sus Principios de Economía (1890) define la economía como la ciencia de la vida diaria en lo que respecta a las acciones humanas tomadas para alcanzar un nivel máximo de bienestar.
Esta definición nos muestra cómo uno de los principios subyacentes a la reflexión económica, pero particularmente enfatizado en la teoría neoclásica, es el del individualismo metodológico . Se contempla el objeto de la teoría como una realidad social compuesta de individuos que se interrelacionan en economías descentralizadas.
En su objetivo fundamental de comprender y predecir el funcionamiento de los mercados, la microeconomía  examina el comportamiento de dos agentes fundamentales: consumidores y productores .
En la teoría de la demanda del consumidor, los individuos objeto de estudio son los consumidores . Se asume que estos se comportan de manera optimizadora y quedan caracterizados por su deseo de consumir ciertos bienes sometidos a una restricción presupuestaria.

En esta exposición, estudiaremos el rol de los hogares en los mercados de bienes:
1.	En este sentido, buscamos modelizar las decisiones del consumidor para explicar el comportamiento de los consumidores en el mercado.
2.	Para ello, revisaremos la teoría de la demanda del consumidor, que busca describir y explicar las elecciones de cestas de consumo observadas.

\subsection{Contextualización}\label{context}


\subsubsection{Relevancia}\label{importance}
Esta cuestión es de gran relevancia a nivel teórico, tanto a nivel microeconómico como a nivel macroeconómico:
•	A nivel microeconómico, supone un paso previo para la determinación del equilibrio de mercado.
•	A nivel macroeconómico, conocer cómo se comportan los consumidores es un paso previo a conocer cómo se comporta el consumo en términos agregados, por lo que el estudio de la demanda del consumidor tiene gran interés si tenemos en cuenta la teoría macroeconómica y la trascendencia del consumo en la demanda agregada.

\subsubsection{Historia}\label{history}

Desde un punto de vista histórico, el estudio de la modelización del comportamiento de los consumidores tiene su origen en el concepto de utilidad , acuñado por autores como BENTHAM, BERNOULLI y GOSSEN. La utilidad se define como la capacidad de los bienes de satisfacer necesidades:
•	Ya en el siglo XVIII, DANIEL BERNOULLI realizó cálculos analíticos y gráficos sobre la utilidad marginal. Este autor realiza el primer análisis riguroso de la toma de decisiones en ausencia de certidumbre, dando una respuesta a la conocida paradoja de San Petersburgo [ver tema 3.A.10] que no parecía resolverse aplicando el criterio de la ganancia esperada.
•	En 1844, JULES DUPUIT relaciona la demanda con la utilidad del consumidor e introduce el concepto de excedente del consumidor.
•	En 1854, HERMANN HEINRICH GOSSEN formuló lo que se conocen como las 3 leyes de Gossen:
	Ley de la saturación de necesidades: “La utilidad de un bien disminuye conforme aumenta la cantidad disponible de ese bien, hasta llegar a la saciedad”. La utilidad marginal es decreciente.
	Ley de igualdad de la utilidad marginal: “Una persona distribuye su ingreso entre los distintos bienes de modo que el grado final de utilidad de cada bien sea el mismo”.
	Ley de la escasez: “Un bien tiene valor únicamente cuando su demanda excede a su oferta”. La escasez es una condición previa para el valor económico. La escasez se refleja en el problema del consumidor en que se enfrenta a unos precios estrictamente positivos, pues sin escasez todo sería gratis (la oferta siempre superaría la demanda).
Más tarde, con la Revolución Marginalista (1870s), autores como LÉON WALRAS (1874) o ALFRED MARSHALL (1890) siguen a GOSSEN y se basan en el análisis marginal, por el cual el consumidor maximiza una función de utilidad sujeta a una restricción presupuestaria y donde se obtiene que el consumidor distribuye su renta entre los distintos bienes hasta que el cociente de las utilidades marginales entre sus propios precios se iguala (es decir, se cumple la segunda ley de Gossen). Este resultado nos permite construir una función de demanda, que relaciona las variaciones en el propio precio con la cantidad deseada por el consumidor.
Estos autores (ALFRED MARSHALL y LÉON WALRAS) recurren a una concepción cardinal de la utilidad, que permitía comparar entre niveles de utilidad y no solamente ordenar las distintas cestas de consumo.
Aunque FRANCIS YSIDRO EDGEWORTH  (1881) ya había descrito la utilidad como función genérica de las cantidades de todos los bienes y había ideado las curvas de indiferencia, serían IRVING FISHER (1892) y VILFREDO PARETO (1896) los primeros en recurrir a una visión ordinal de la utilidad  y en reconocer que cualquier transformación creciente arbitraria de la función de utilidad no tenía efecto alguno sobre la demanda. Hoy en día se acepta ampliamente en teoría de la demanda que solo la utilidad ordinal importa. Es decir, una función de utilidad sirve meramente como instrumento conveniente para representar una relación de preferencia, y cualquier transformación creciente de la misma servirá también a este propósito.
Sin embargo, no fue hasta los años 30 cuando JOHN HICKS y ROY ALLEN (1934) introducen el concepto de preferencias, que daba generalidad a la teoría.
Cabe mencionar además la contribución de GÉRARD DEBREU (1959) en cuanto a la axiomática que dota de una mayor formalidad a esta teoría.
Por otro lado, el ingeniero italiano GIOVANNI B. ANTONELLI (1886) daría origen a lo que hoy se conoce como teoría de la dualidad al hacer el camino en sentido contrario: construir curvas de indiferencia y una función de utilidad a partir de funciones de demanda inversas. Para esta recuperación de la función de utilidad se han seguido dos vías:
a)	Por un lado, la vía de la preferencia revelada de PAUL A. SAMUELSON (1947) y HENDRIK S. HOUTHAKKER (1950); y
b)	Por otro lado, la resolución del denominado problema de la Integrabilidad de la mano de LEONID HURWICZ y HIROFUMI UZAWA (1971), buscando ciertas propiedades de las funciones de demanda como las referidas a la matriz de Slutsky.
Serían KIHLSTROM, MAS-COLELL y SONNENSCHEIN (1976) quienes unificarían ambos enfoques relacionando los axiomas de la preferencia revelada con las propiedades de la matriz de Slutsky.
Para el desarrollo de la teoría de la demanda del consumidor el análisis de dualidad ha sido determinante. En origen, el análisis de dualidad se pregunta: si “se da la vuelta” al problema de maximización de preferencias, de modo que el consumidor persiga minimizar el gasto en bienes restringido por alcanzar un nivel de utilidad determinado, ¿cambiarían las decisiones del consumidor? 
Lo que en apariencia es una cuestión teórica casi anecdótica se revela de enorme utilidad, pues no solo consigue responder a esta cuestión, sino también, por el camino proveer de una serie de relaciones y propiedades del problema del consumidor que enriquecen enormemente la teoría neoclásica de la demanda de consumo, permitiendo contestar a preguntas como:
1.	¿Qué demandas observadas son compatibles con la maximización de preferencias?
2.	Dadas unas condiciones necesarias para esta compatibilidad, ¿son también suficientes?
3.	¿Qué tipo de preferencias genera una clase particular de demandas?
Como se puede apreciar, la moderna teoría de la demanda busca ligar ambos conceptos: preferencias y demanda. Y el análisis de dualidad ha sido determinante para resolver estas cuestiones.
Antes de comenzar con la exposición, se hace necesario delimitar el significado del término “dual”. En principio, “dual” puede entenderse en un sentido estricto para referirse exclusivamente a la relación entre el problema de maximización de la utilidad y el de minimización del gasto, en que se invierten los papeles de función objetivo y restricción. No obstante, el término “dual” se suele emplear en la actualidad en economía en un sentido más amplio, aplicado a cualquier par de problemas o incluso conceptos que son formalmente similares excepto por el rol invertido, no solo de función objetivo y restricción, sino también de cantidades y precios, o de maximización y minimización (MAS-COLELL, WHINSTON y GREEN, 1995; BLUME, 2008). A lo largo del tema, el alcance del término “dual” será el de su versión restringida, aunque no debe extrañar que se utilice en algún momento en sentido amplio.\footnote{%
    Por ejemplo, en la relación de dualidad entre función indirecta de utilidad y función de utilidad. En la teoría de costes y del beneficio, esta versión amplia del término “dual” es más habitual, por ejemplo cuando se habla de la dualidad entre el problema de minimización del coste y maximización del beneficio (no son problemas “duales” en sentido estricto pues el problema del beneficio incorpora el elemento “ingresos” del que carece el problema del coste) o la dualidad entre las funciones de coste/beneficio y la tecnología subyacente (aludiendo a la recuperación de la última a partir de las primeras) [ver tema 3.A.12].
}
La teoría de la dualidad, en su aplicación a la teoría de demanda del consumidor, es sin duda uno de los desarrollos más importantes (si no el más importante) en el ámbito de la teoría del consumidor:
	Por un lado, refuerza las conclusiones principales de dicha teoría, dotando de rigurosidad a la teoría de la demanda del consumidor.
	Por otro lado, ofrece una solución para superar su principal limitación: el hecho de que la función de utilidad, necesaria para obtener la función de demanda, no sea observable.
Además, ofrece aplicaciones muy relevantes para ciertos campos, como por ejemplo la Economía del Bienestar.


\subsubsection{Problemática}\label{problem} 
1.	¿Qué demandas observadas son compatibles con la maximización de preferencias?
2.	Dadas unas condiciones necesarias para esta compatibilidad, ¿son también suficientes?
3.	¿Qué tipo de preferencias genera una clase particular de demandas?




\subsection{Estructura de la exposición}\label{structure}
1.	Análisis teórico: Especificación y resolución de los problemas primal y dual y relaciones de dualidad entre ambos problemas y sus resultados.
a.	Esta parte, un tanto teórica y abstracta, tiene su razón de ser en el siguiente apartado.
2.	Aplicaciones prácticas: Utilidad práctica de este análisis de dualidad en el campo de la teoría de la demanda de consumo neoclásica:
a.	El problema de la Integrabilidad de las preferencias;
b.	La ecuación de Slutsky; y
c.	El análisis de bienestar.




\clearpage
\section{Primal vs. Dual}

\imagenfit{Primal_Dual.png}{Esquema de relaciones entre el problema primal y su dual}{Apuntes Víctor Gutíerrez. Tema 3.B.9}

\subsection{Teoría de la Dualidad: Nociones básicas y definiciones}
Para modelizar el comportamiento de un consumidor podemos adoptar 2 enfoques:
a)	El enfoque de la utilidad: consiste en buscar la cesta de consumo que maximicen la utilidad del consumidor, sujeto a una restricción de recursos o presupuestaria (problema de maximización de la utilidad). Por tanto, los recursos se toman como dados (son una restricción) y se busca la mejor manera de combinar los diferentes bienes disponibles para alcanzar la mayor utilidad posible (es la función objetivo); o,
b)	El enfoque del gasto: que consistiría en la inversión de este mismo, de modo que se pasa a considerar como fijo o dado un nivel concreto de utilidad, que se convierte en la restricción, y se busca la combinación de bienes que permita alcanzarlo incurriendo en el menor gasto posible, que deviene la función objetivo (problema de minimización del gasto).
El desarrollo formal de este marco fue obra de MAS-COLELL, WHINSTON y GREEN (1995) que, entre otras, gracias a las aportaciones al campo al campo de la programación lineal de VON-NEUMANN (1947), unificaron los dos enfoques mencionados en un mismo marco, tratándose de esta manera de una “Relación DUAL” que presenta unas relaciones características, que son las que vamos a pasar a estudiar. 
Nos referiremos al Problema de Maximización (Max) como Problema Primal y al Problema de Minimización como su “Dual”. Principalmente porque el Problema de Maximización fue objeto de desarrollo en el Tema 8 [ver Tema 3.A.8] por lo que nos limitaremos a plantear su dual y las relaciones que subyacen.

\subsubsection{Marco común a los dos porblemas}
El enfoque dual requiere exactamente la misma especificación del conjunto de elección, preferencias y conjunto presupuestario que el enfoque primal [del que aquí se dará un resumen por ser objeto de un mayor desarrollo en el Tema 3.A.8]:
1.	Partimos de un conjunto de elección no vacío definido en el espacio R^K+, compuesto por subconjuntos o cestas de consumo sobre las que el agente tiene unas preferencias racionales.
2.	Estas preferencias, que son relaciones binarias que permiten al consumidor comparar las diferentes cestas de consumo, dan lugar a una ordenación completa y transitiva de las diferentes cestas de consumo.
Aunque aquí acaba la definición axiomática “libre de críticas”, por decirlo de alguna manera, añadiremos, para facilitar nuestro análisis tres restricciones (o axiomas) adicionales:
3.	Deseabilidad. Monotonía en sentido fuerte.
4.	Convexidad. De las preferencias.
5.	Diferenciabilidad.  
Estos axiomas garantizan que las preferencias sean de “buen comportamiento” y podremos hallar una Función de Utilidad del mismo tipo (“buen comportamiento”, cuasi-cóncava y diferenciable), imprescindible para plantear ambos de optimización.
Además, es relevante tener en cuenta que el individuo hace frente a un gasto, derivado de que los precios de los bienes son positivos, que:
a)	En el Problema Primal esto dará lugar a una restricción presupuestaria; y,
b)	En su “Dual” el gasto constituirá la función objetivo a minimizar.

\subsection{Problema Primal. Maximización de la Utilidad}
El enfoque primal de la demanda de consumo, y por ende el enfoque dual, busca resolver un problema económico, es decir, conjugar unas necesidades a priori infinitas con unos recursos escasos. Para ello, tras haber especificado ambos componentes (preferencias y recursos), se hace a continuación un breve repaso del análisis primal.

\subsubsection{Especificación}
Dicho esto, recordamos que el problema primal del consumidor se especificaría mediante el siguiente problema de optimización condicionada:

\eqblock{
\begin{aligned}
    \boxed{\begin{aligned}
        &\max{U(x)}\\[0.5em]
        &\text{s.a.}\quad 
        \begin{cases}
            p\cdot x\leq\bar W\\[0.5em]
            0\leq x
        \end{cases}
    \end{aligned}}
\end{aligned}
}{}

Con una función de utilidad cuasi-cóncava (podemos suponer que será estríctamente cuasi-cóncava si imponemos preferencias estríctamente convexas) y un conjunto presupuestario definido, no vacío, cerrado y acotado (inferior y superiormente). Con todo esto, se puede asegurar que el óptimo existe, es global y es único.

\subsubsection{Desarrollo formal}
Si restringimos la economía a dos bienes, podemos dar un ejemplo gráfico muy simple de todas las implicaciones que hemos mencionado. 

[INCLUIR GRÁFICOS QUE ILUSTREN GRÁFICAMENTE EL DESARROLLO DEL PROBLEMA]

\subsubsection{Solución: SCDEM}
Como hemos visto la solución al problema de maximización de la utilidad es una cesta de bienes que permite obtener la mayor utilidad posible dados unos precios y un presupuesto. 
Sin embargo, existe también la posibilidad de formular la ecuación como sistema, gracias a las características del sistema de ecuaciones que forman el conjunto de CPO’s, se puede formular en forma de función, dependiente de los precios y el nivel de renta, de tal forma que no debamos recurrir al problema de maximización cada vez que debamos resolverlo.\footnote{%
    Esto nos permitirá realizar análisis de estática comparativa y estudiar cómo se ven afectadas las demandas del consumidor ante cambios en los parámetros del problema.
} Esto es lo que se conoce como Sistema Completa de Ecuaciones de Demanda Marshalliana y, no es más que la formulación del problema de maximización como el del argumento del máximo, esto es, el argumento del máximo de la Función de Utilidad, dadas las restricciones presupuestaria (i.e. producto vectorial igual a la dotación) y de no-negatividad. 
x(p,M)=(■(x_1^*@⋮@x_i^* ))=(■(x_1 (p_1,p_2,⋯,p_i,i)@⋮@x_i (p_1,p_2,⋯,p_i,i)))
Presenta las siguientes propiedades:
	El SCEDM existe, es global y es único (Teorema de la Función Implícita).
	Las funciones de demanda son continuas en tanto en precio como en renta (Teorema del Máximo).
	Las funciones de demanda son homogéneas de grado cero tanto en precios como en renta, lo que se traduce en que el agente no adolece de ilusión monetaria, es decir, ante aumentos en la misma proporción de precios y riqueza las decisiones de consumo no se ven alteradas.\footnote{%
	    Si se multiplican todos los precios y la renta por una constante, las restricción presupuestaria queda inalterada, luego el problema y su solución también quedan inalterados.
	} Esto permite prescindir de la referencia al dinero (y a los precios absolutos o nominales) en cualquier análisis de demanda, simplemente designando a cualquiera de los n bienes como “numerario” o dinero y dividiendo todos los precios nominales y la riqueza entre el precio del numerario, de modo que la demanda dependa exclusivamente de la riqueza y los precios relativos o reales.
	Se cumple la Ley de Walras, esto es, en el óptimo la restricción presupuestaria siempre se satura (se cumple con igualdad estricta), lo que se traduce en que el consumidor siempre agota su riqueza disponible, pues no sería racional dejar riqueza sin gastar debido a las implicaciones del axioma de deseabilidad.\footnote{%
	    Se le denomina así por la Ley de Walras que se obtiene en equilibrio general, p ∙ z(p) = 0, cuya demostración matemática revela que este producto de precios por excesos de demanda no es sino la agregación de las restricciones presupuestarias de todos los consumidores de la economía, y en el agregado también deben agotarse todos los recursos disponibles [ver tema 3.A.21].
	}
El SCEDM es, con cierta generalidad, diferenciable, tanto respecto a precios como respecto a renta.\footnote{%
    Los requisitos para que el SCEDM sea diferenciable (JEHLE y RENY) si asumimos solución interior (cantidades estrictamente positivas de todos los bienes) y precios y riqueza estrictamente positivos son:
1.	Función de Utilidad continua y dos veces diferenciable;
2.	∂U(x*) ∂xi ⁄ > 0 para algún i ∈ (1, … ,n);
3.	La matriz Hessiana de la Función de Utilidad tiene un determinante distinto de 0 en el óptimo.
}

\subsubsection{Función Valor: Función Indirecta de Utilidad}
Dicho esto, veremos la primera herramienta que necesitaremos para establecer nuestras relaciones de dualidad. 
Un Problema Primal (de éstas características) nos permite obtener una herramienta adicional, que dará a conocer los valores que la función objeto de optimización (i.e. la Función de Utilidad) toma cuando se alcanza el equilibrio (esto es, cuando dados unos precios y un nivel de riqueza, se escoge/encuentrala cesta de consumo óptima).
Este herramienta, conocida como Función Valor o Función Indirecta de Utilidad (pues nos proporciona sus niveles pero indirectamente), se obtiene sustituyendo la Demanda Óptima del problema (o conjunto de Demandas Óptimas), en la función objetivo (la Función de Utilidad).
[DEFINIR MATEMÁTICAMENTE LA FUNCIÓN VALOR / FIU]
Por tanto, la FIU permite conocer los niveles de utilidad óptimos para cada conjunto dado de parámetros (precios y riqueza), sin necesidad de repetir la resolución del problema primal. De ahí su nombre y su interpretación: se obtiene la utilidad del consumidor a través de los “resultados” que daría el problema de maximización a cuyas restricciones está sujeto, frente a la Función de Utilidad (convencional) que permite obtener la Utilidad que el consumidor percibe de cada cesta (o combinación) de consumo.
(Atendiendo a las especificaciones dadas) Podemos comentar algunas de las propiedades de la FIU que más adelante serán de utilidad:
1.	La FIU existe;
2.	La FIU es continua en precios y renta;
3.	La FIU es homogénea de grado cero tanto en precios como en renta. Es decir, un cambio equiproporcional en todos los precios y riqueza deja inalterado el equilibrio del consumidor (debido a la homogeneidad de grado cero de las funciones de demanda) y, por lo tanto, también su bienestar.
4.	Es estrictamente creciente en la riqueza. Es decir, una mejora en el presupuesto del agente debido a un aumento de la riqueza aumenta siempre el bienestar.
5.	Es no-creciente (es decir, decreciente pero no estrictamente) en los precios. Es decir, un empeoramiento en el presupuesto del agente por un aumento de los precios nunca podrá aumentar el bienestar. Como mucho podrá mantenerse constante ante una solución de esquina.
6.	Es cuasi-convexa tanto en precios como en renta. Es decir, debido a la convexidad de las curvas de indiferencia la utilidad máxima que puede alcanzarse con presupuestos promedios (aquéllos donde los precios de los bienes se distribuyen de manera “relativamente homogenea”) es inferior a la que puede alcanzarse con presupuestos extremos (JEHLE y RENY, 2011), i.e. el consumidor alcanza un mayor nivel deutilidad en casos donde los bienes guarden una relación de precios “desproporcionada”.
7.	Cumple la identidad de Roy, si la FIU es diferenciable, entonces: 
 
En palabras, las demandas ordinarias pueden encontrarse directamente a partir de la FIU (mediante derivación parcial) [para la demostración ver anexo A.1].
Ahora, cabría preguntarse: Partiendo desde el mismo contexto axiomático, liberando la condición previa de optimización dada por lo parámetros: ¿se alcanzaría, para un mismo nivel de utildad, el equilibrio en los mismos términos? 
Responder a esta inocente pregunta nos permitirá descubrir una serie de propiedades y relaciones entre funciones de ambos enfoques (primal y dual) que enriquecen enormemente la teoría de la demanda de consumo neoclásica.

\subsection{Problema Dual: Minimización del gasto}
Vamos a adentrarnos en el problema de optimización alternativo que nos propone la Teoría de la Dualidad, el de “minimización del gasto”.

\subsubsection{Especificación}
El PMG se plantea entonces como un problema de optimización condicionada que especificaemos como sigue (recuérdese que se parte de los mismos supuestos o axiomas que en el problema anterior, mencionados al principio):

\eqblock{
\begin{aligned}
    \boxed{
    \begin{aligned}
        \min_{h} \quad p\cdot h\\[0.5em]
        \text{s.a.}\quad
        \begin{cases}
            u(h)\ge\bar U\\[0.5em]
            h\ge 0
        \end{cases}
    \end{aligned}
    }
\end{aligned}
}
 
Se trata de minimizar el gasto o desembolso del consumidor dados unos precios, fijando como restricción un nivel de utilidad mínimo a obtener. En otras palabras, se busca alcanzar (al menos) el nivel de utilidad  de la forma más barato posible dados unos precios.

\subsubsection{Desarrollo}
Dado que se trata de un problema de optimización condicionada por restricciones de desigualdad, formularemos el Lagrangiano:

\eqblock{
\begin{aligned}
    \mathcal{L}=p\cdot h+\mu\cdot(\bar U-u(h))
\end{aligned}
}

Y, al evaluar las CPO’s, deberemos comprobar si se cumplen las condiciones necesarias (KHUN-TUCKER) para dar lugar a una solución (alternativamente, podemos también identificar si se dan las condiciones suficientes para establecer que se trata de una solución al probema de optimización, si se dan las de primer orden, CPO’s, se dan también las de seguundo orden, CSO’s). Sin entrar en el desarrollo analítico de la ecuación para ningun escenario definido simplemente cabedestacar que, diremos que es solución (único vector) al problema siempre y cuando cumpla las siguientes condiciones:
1.	Condición de estacionariedad. En caso de solución interior, los precios ponderados por las utilidades marginales se igualan entre sí para los respectivos bienes y a su vez a mu.
2.	Restricción de utilidad. La restricción se satura y el consumidor alcanza exactamente la utilidad requerida, fruto del axioma de monotonía de las preferencias.
3.	No negatividad de los bienes. En el óptimo el consumidor adquiere una cantidad positiva o nula de los bienes.
4.	Condición de holgura complementaria: $h_i^*\cdot [p_i-\mu\cdot u_i(h)]=0$
 
Se debe anular la condición (1) (solución interior), la condición (3) (solución de esquina), o ambas a la vez, tal que se de una solución de esquina. De ahí que el producto de ambas siempre sea nulo.
[MOSTRAR GRÁFICAMENTE LO QUE ESTO SIGNIFICA]


\subsubsection{Solución: SCEDC}
Por último, cabe destacar que, aunque la solución a este problema de minimización del gasto es la cesta de bienes más barata que permite alcanzar un nivel de utilidad dado. 
Si levantamos la restricción de estricta convexidad, que nos ha faciliitado su representación, esta solución podría dar lugar a lo que formalmente se conoce como Sistema Completo de Ecuaciones de Demanda Compensada (o Hicksiana), que se trata de un vector de cantidades demandadas óptimas en función de un nivel de precios y utilidad dados:

\eqblock{
\begin{aligned}
    h^*=h(p,\bar U)=
    \left(\begin{matrix}
        h_1^*\\
        \vdots\\
        h_1^*
    \end{matrix}\right)=
    \left(\begin{matrix}
        h_1(p_1,p_2, \dots ,p_n,\bar U)\\
        \vdots\\
        h_n(p_1,p_2, \dots ,p_n,\bar U)
    \end{matrix}\right)
\end{aligned}
}
 
El calificativo de compensadas proviene del hecho de que el presupuesto o gasto ya no es fijo si no variable, de modo que ante cualquier variación en los precios, la utilidad se mantendrá constante y dicho gasto se reajustará para acomodar o compensar dicha variación.
El SCEDC posee las siguientes propiedades:
1.	El SCEDC existe, es global y es único;
2.	Las funciones de demanda son continuas, tanto en precios como en Utilidad;
3.	Las funciones de demanda son homogeneras de grado cero en precios, por lo que se traduce en que el agente no adolece de ilusión monetaria y, por tanto, si todos los precios se multiplican por una constantepositiva, las demandas hicksianas no varían:
$$h(\theta\cdot p,\bar U)=\theta^0\cdot h(p,\bar U)=h(p,\bar U)$$
4.	Las funciones de demanda compensadas cumplen la Ley de la demanda, esto es, precio y cantidad nunca se moverán en la misma dirección (no hay posibilidad de bienes Giffen en las demandas hicksianas).\footnote{%
    Recuérdese que la razón de ser de un bien Giffen es el efecto renta (que es mayor en valor absoluto al efecto sustitución y en sentido contrario), y este efecto está ausente en la demanda hicksiana. Esta propiedad es clave y es la diferencia principal con la demanda Marshalliana. Esto se extrae de que la matriz de Slutsky (matriz de primera derivadas) es semidefinida negativa.
} Esta propiedades clave y es la diferencia fundamental con la demanda marshalliana. Recordemos esta propiedad pues extraeremos algunas conclusiones muy importantes de ella más adelante.

\subsubsection{Función Valor: Función de Gasto}
De la misma forma que en el PMU, el problema dual permite obtener una función a través de la cual conocer los valores que la función objeto de minimización toma cuando se alcanza el equilibrio. Es decir, obtener el valor de la restricción (frontera superior) dados unos precios y un nivel de utilidad requerido
Formalmente se conoce como Función Valor (óptimo), o Función de Gasto (para este caso), del problema, que se obtiene sustituyendo la ecuación (o ecuaciones, en el caso generalizado: levantamos restricciones) de demanda óptima en la función objeto de minimización:\footnote{%
    Nótese que: La Función de Gasto recupera el mínimo de la función objetivo (min); mientras que el SCEDC recupera la cesta óptima asociado al gasto mínimo, lo que se conoce como argmin
}

\eqblock{
\begin{aligned}
    E(p,\bar U)=p\cdot h(p, \bar U)=\min_{h}\{p\cdot h\;|\;u(h)=\bar U,h\ge0\}
\end{aligned}
}

Es decir, la FG permite obtener los niveles de gasto óptimos (mínimos) para cada conjunto dado de parámetros (precios y utilidad mínima), sin necesidad de volver a resolver el problema dual cada vez que se modifiquen estos parámetros.
La FG posee unas propiedades de interés, entre las que cabe destacar:
1.	Su continuidad, tanto en precios como en renta;
2.	Su homogeneidad de grado UNO en precios. Es decir, un cambio equiproporcional en todos los precios, aunque deja inalterado el equilibrio del consumidor (debido a la homogeneidad de grado cero de las funciones de demanda), sí afecta de manera equiproporcional al gasto incurrido;
3.	Es estríctamente creciente en la Utilidad;\footnote{%
    Esto es así porque las curvas de indiferencia nunca se cortan (levantemos o no los supuestos restrictivos iniciales).
}
4.	Es no-decreciente en precios (es decir, creciente, pero no estrictamente [siempre que levantemos las condiciones impuestas al principio]\footnote{%
    Si levantamos las restricciones de estricta convexidad de las preferencias, tendremos una función de utilidad cuasi-cóncava (no estríctamente cuasi-cóncava), por lo tanto, podría ser que, ante un cambio de precios (por la saturación en el equilibrio de diferentes cestas de demanda óptima), la utilidad permanezca inalterada al distribuir la proporción de bienes en las cestas de consumo óptima [incluso podrá dar lugar a una única cesta de consumo óptima en el mismo nivel de equilibrio (en términos de utilidad) donde antes se tenían varias].
})
Además, ciertas propiedades llevan a reconocer ciertas propiedades de la denominada “matriz de Slutsky” o “matriz de sustitución”: S.\footnote{%
    La matriz Hessiana o de Segundas Derivadas, recoge las segundas derivadas de una Función. En este caso, la matriz de Slutsky es la matriz de segundas derivadas de la Función de Gasto respecto a todos los precios, es decir, la matriz que recoge todos los efectos precio, tanto propios como cruzados, sobre las demandas hicksianas.
}
1.	Su concavidad en precios (para un nivel de Utilidad dado) [levantando los supuestos de estricta-convexidad]. Es decir, Ante aumentos en los precios, el consumidor podría comportarse de manera “pasiva” y no modificar la cesta de consumo original, manteniendo la cesta original, pero como hemos visto el gasto también aumentará. Por tanto, un comportamiento optimizador lo llevará a sustituir unos bienes por otros para reducir el gasto necesario para alcanzar dicha utilidad inicial. En resumen, dicho de manera intuitiva, aumentos en los precios llevan a aumentos en el gasto menos que proporcionalmente (de ahí la concavidad), gracias a estas posibilidades de sustitución entre bienes.
2.	El cumplimiento del Lema de Shepard. Esto es, si la Función es diferenciable, entonces las demandas hicksianas pueden encontrarse directamente a partir de la Función de Gasto (mediante derivación parcial) [demostración en el Anexo A.2]

$$\begin{aligned}
    S=
    \left(\begin{matrix}
        s_{11}&\cdots&s_{1n}\\
        \vdots&\ddots&\vdots\\
        s_{n1}&\cdots&s_{nn}
    \end{matrix}\right)\quad \text{donde,}\;\;
    \left\{\begin{aligned}
        s_{ij}\equiv\frac{\partial^2E(p,\bar U)}{\partial p_ip_j}\underset{
        \substack{\text{\scriptsize Por el Lema}\\
        \text{\scriptsize de SHEPARD}}{=}\frac{\partial h_i(p, \bar U)}{\partial p_j}\\[0.5em]
        \forall i,j=1,2,\dots,n
    \end{aligned}\right.
\end{aligned}$$
 
Esta matriz es:
1.	Simétrica. De modo que el orden de derivación de la Función de Gasto no altera su significado económico (pues no altera la matriz final), es decir, a diferencia de las demandas marshallianas, los efectos precio cruzados para dos bienes cualesquiera en las demandas hicksianas son idénticos;\footnote{Por el teorema de Young o Lema de Schwarz.
} y,
2.	Semidefinida negativa (por la concavidad de la FG): Lo que implica que, a diferencia de las demandas marshallianas, los efectos precio-propio en las demandas hicksianas son no positivos.\footnote{%
    Esto es así poque todos los elementos de la diagonal principal de una matriz semidefinida negativa deben ser no positivos (nulos o negativos). Y estos elementos en la matriz de Slutsky, S, son los efectos precio propios. De modo que si los efectos precio propios son no-positivos: el aumento de precio de un bien, o bien reduce su demanda Hicksiana o la mantiene constante, pero en ningún caso aumenta (Ley de la Demanda Compensada).
} De ahí que las demandas hicksianas sí cumplan la “ley de la demanda compensada”.

\subsection{Relaciones de Dualidad}
Recapitulando, se han obtenido todos los elementos necesarios (Funciones de Demanda, Funciones de Valor y sus propiedades) para pasar a continuación a establecer las relaciones existentes entre los diferentes componentes de la Relación de “Dualidad”.

\subsubsection{Teorema Básico de la Dualidad: \textit{Dualidad entre soluciones}}
Por el “Teorema Básico de Dualidad” sabemos en primer lugar que las cestas de consumo óptimas de ambos problemas son idénticas siempre que se realice el debido intercambio de información entre problemas: esto es, siempre que se desarrollen dentro del mismo marco normado o axiomático.

[MOSTRAR GRÁFICAMENTE LA RELACIÓN DE DUALIDAD]
A mode de ejemplo, supongamos que tenemos una Función de Utilidad continua que muestra unas preferencias que cumplen la axiomática inicialmente planteada (racionalidad, deseabilidad y convexidad estricta). Entonces, si asumimos que los precios son estríctamente positivos tendremos que la solución óptima (cesta de consumo óptima) del PMU, dada una riqueza estríctamente positiva, condicionaría también la misma solución óptima en el PMG si la Utilidad requerida es precisamente la proporcionada por esta cesta de consumo óptima. Entonces, la Función de Gasto recuperaría precisamente la cota superior del conjunto de gasto que optimiza (minimiza) el PMG.
Es decir, la solución numérica o cesta de cantidades consumidas del problema primal y del problema dual coincidirán si se realiza el adecuado trasvase de información acerca de los niveles de utilidad y de renta. Este resultado ratifica la intuición de equivalencia que se puede demostrar gráficamente de manera directa:

\imagenfit{Dualidad_Soluciones.png}{Representación gráfica del problema primal y el problema dual}{Apuntes Victor Gutierrez. Tema 3.A.9}
 
Ahora bien, este resultado se refiere exclusivamente a las cestas óptimas, pues es evidente que las funciones de Demanda Marshallianas y Hicksianas serán cosas distintas (ya que unas y otras dependen de distintas variables)
No obstante, existe un par de relaciones de equivalencia/dualidad entre dichas funciones que se derivan del “Teorema Básico de Dualidad”. Estas son referidas tanto a las funciones de Demanda (marshallianas y hicksianas) como a las Funciones Valor (FIU y FG).\footnote{%
    En esta aproximación, llamaremos curvas isogasto a las diferentes restricciones que se proyectan sobre el plano, mientras que la curva de indiferencia representa la curva correspondiente al nivel dado de Utilidad.
}

\imagenfit{ProblemaDual.png}{Representación gráfica del problema dual}{Apuntes Víctor Gutierrez. Tema 3.A.9}
  
Un punto interesante sobre la relación de dualidad eque podemos observar gráficamente es la relación inversa que muestran:
a)	Mu ($\mu$) [el multiplicador de Lagrange, en el problema de optimización], y nos informa sobre la sensibilidad de la Función objetivo (de Gasto) ante cambios en la restricción (i.e. la utilidad requerida). En concreto, representa lo que varía el gasto al variar marginalmente la utilidad requerida; y
b)	El multiplicador de Lagrange del problema primal ($\lambda$) (precio sombra o utilidad marginal de la riqueza, [ver Tema 3.A.8]
Esto significa que cuando el multiplicador de Lagrange del problema primal aumenta, el multiplicador de Lagrange del problema dual disminuye, y viceversa. 

Además, para el caso de dos bienes tendríamos que, la RMS se iguala al cociente de precios. Esto es, se cumple la segunda Ley de Gossen: “La igualdad de las utilidades marginales por sus propios precios”. Esto implica que el consumidor decidiría su demanda igualando la tasa objetiva de intercambio de los dos bienes en el mercado (esto es, los precios) con la tasa subjetiva que depende de sus preferencias (utilidad marginal):

$$\begin{aligned}
    \text{RMS}_{x_1}^{x_2}=-\frac{UMg_1}{UMg_2}=-\frac{p_1}{p_2}
\end{aligned}$$

\subsubsection{Dualidad entre Funciones de Valor}
La Función Indirecta de Utilidad es la inversa de la Función de Gasto (o viceversa) bajo unas condiciones muy generales. En efecto, hallada una de las funciones valor, mediante una sustitución y una sencilla reorganización de la ecuación, puede alcanzarse la otra sin necesidad de resolver su problema correspondiente (hallada la FIU se puede encontrar directamente la FG sin necesidad de resolver el problema dual, o viceversa):\footnote{Por el teorema de Young o Lema de Schwarz.
}

$$\begin{aligned}
    E(p,\bar U)\;\overset{\bar U=V(p,w)}{\Longrightarrow}\;E(p,V(p,\bar W))=\bar W\;\overset{V=E^{-1}}{\Longrightarrow}\;V(p, \bar W)=E^{-1}(p,\bar W)\\[0.5em]
    V(p,\bar W)\;\overset{\bar W=E(p,\bar U)}{\Longrightarrow}\;V(p,E(p,\bar U))=\bar U\;\overset{E=V^{-1}}{\Longrightarrow}\;E(p, \bar W)=V^{-1}(p,\bar U)
\end{aligned}$$
  
Matemáticamente, esto responde a que los conjuntos sobre los que se evalúa cada Función tienen una relación de biyectividad. Es decir, (en términos económicos) que para todo nivel de utilidad óptimo (máximo) hay asociado un nivel de gasto óptimo (mínimo), y viceversa [Si dichos conjuntos no fueran biyectivos, no podríamos establecer esta relación].




\clearpage
\section{Análisis Dual: Aplicaciones}
El análisis de dualidad va más allá de establecer una una relaciones teóricas entre problemas primal y dual, pues tiene relevantes aplicaciones prácticas:
1.	La dualidad permite establecer formalmente la ecuación de slutsky, descomponiendo el efecto total de una variación de precios en efecto, sustitución y efecto renta, ya sea propia o cruzado.
2.	Además, la actualidad permite resolver la cuestión de la integralidad, esto es, recuperar las preferencias (no observables) a partir de demandas (observables). Esto, a su vez, tiene aplicaciones prácticas a la hora de estimar a partir de datos reales sistemas de ecuaciones de demanda.
Finalmente, la dualidad permite desarrollar herramientas que cuantifiquen el cambio en el bienestar ante cambios de los parámetros, precios y renta, tanto con información completa (variación equivalente y conversatorio) como incompleta aproximación mediante la variación del excedente del consumidor e índice de precios.

\subsection{Ecuación de SLUTSKY}
La dualidad entre funciones de demanda permite obtener una derivación formal de la ecuación de Slutsky, según la cual es posible descomponer el efecto sobre una demanda marshalliana de la variación en un precio (propio o cruzado).

\subsubsection{Derivación formal}
Por el teorema de la dualidad sabemos que en el óptimo las demandas son idénticas (tanto de manera compuesta, la cesta de bienes completa, como a título particular para cada Bien i). 
1.	Por tanto, tenemos que:

$$\begin{aligned}
    h_i(p,\bar U)=x_i(p,E(p,\bar U))
\end{aligned}$$

2.	[Nótese que estamos trabajando sobre la demanda de un Bien i particular, aunque la igualdad se mantiene puesto que son idénticas las demandas]. A partir de esta igualdad, podemos derivar respecto al precio de un Bien j obteniendo:

$$\begin{aligned}
    \frac{\partial h_i(p,\bar U)}{\partial p_j}=\frac{\partial x_i(p,E(p,\bar U))}{\partial p_j}\\[0.5em]
    \frac{\partial h_i(p,\bar U)}{\partial p_j}=\frac{\partial x_i(p,\bar W)}{\partial p_j}+\frac{\partial h_i(p,\bar W)}{\partial E(p,\bar U)}\cdot \frac{\partial E(p,\bar U)}{\partial p_j}
\end{aligned}$$
 
3.	[Por la regla de la cadena: parte dcha. de la ecuación] A partir de la ecuación anterior, podemos hacer uso de las relaciones de dualidad, de modo que:

$$E(p, \bar U)=\bar W$$
 
4.	Y, sabiendo que, por el Lema de Shepard:\footnote{%
    El lema de Shephard afirma que la demanda de un bien particular para un determinado nivel de utilidad, dados los precios p, es igual a la derivada de la función de gasto con respecto al precio del bien Tcorrespondiente
}
$$\frac{\partial E(p,\bar U)}{\partial p_j}=h_j(p, \bar U)=h_j(p,V(p, \bar W))=x_j(p, \bar W)$$

5.	Tenemos que:

\eqblock{
\begin{aligned}
    \underbrace{\frac{\partial x_i(p,\bar W)}{\partial p_j}}_{\text{\scriptsize Efecto Total}}=\underbrace{\frac{\partial h_i(p,\bar U)}{\partial p_j}}_{\text{\scriptsize Efecto sustitución}}-\underbrace{x_j\cdot \frac{\partial x_i(p,\bar W)}{\partial \bar W}}_{\text{\scriptsize Efecto renta}},\quad \forall i,j
    
\end{aligned}
}



 












\clearpage
\section{Conclusión} 

\subsection{Extensiones y relación con otras partes del Temario}


\subsection{Opinión}



\subsubsection{Idea final (Salida o cierra)}


\clearpage
\pagenumbering{gobble}
\MyBibliography
\MyBibItem{Autor}{Título}{Año}{DOI -- LINK}


%ANEXOS
\clearpage
\anexo{\textbf{Preguntas Test}}
%\input{\TemaCode/questions.tex} 




\end{document}